\chapter{Teoretická část}

\section{Funkce expandérů provozního kondenzátů (Flash Tank)}
Flash tanks are widely used across various industrial sectors wherever high-pressure liquids or condensate streams undergo pressure reduction. In power engineering, they are commonly applied in steam turbine systems for handling drains and auxiliary steam flows. In the chemical and petrochemical industries, flash tanks are used for phase separation following throttling processes and pressure letdown operations. In the food and process industries, they are employed in thermal processes such as evaporation, sterilization, and heat recovery systems, where controlled flashing and separation are required. Despite differences in operating conditions and media, the underlying principle remains the same: rapid pressure reduction leading to phase change, followed by separation of vapor and liquid phases.

In general, flash tanks are used across various industrial sectors. However, this work focuses specifically on flash tanks employed in steam turbine systems, where they process turbine drains and auxiliary steam flows and are directly connected to the condenser via piping. These systems operate under conditions typical for power plant applications, including low-pressure environments (often near vacuum) and high volumetric flow rates of the vapor phase. The analysis presented in this thesis is therefore limited to this class of flash tanks, with emphasis on their role within the steam cycle and their interaction with downstream components such as the condenser.

When condensate at elevated pressure is discharged to a lower-pressure environment, it contains excess thermal energy relative to the saturation state at the new pressure. As a result, a portion of the liquid instantaneously evaporates, forming so-called flash steam. This process, known as flash evaporation, leads to the formation of a two-phase mixture consisting of vapor and liquid \cite{wang_construction_2023}. The presence of this mixture requires subsequent phase separation to ensure correct operation of downstream components.

The primary function of the flash tank is therefore twofold: (i) to provide a controlled environment for pressure reduction and phase change, and (ii) to enable effective separation of steam and liquid phases. The separated steam is discharged through a vapor outlet and routed via piping to downstream equipment, most commonly to a condenser. The liquid phase accumulates at the bottom of the vessel and is removed either by means of a siphon or by pumping, depending on the specific design configuration. The detailed implementation of these discharge methods is discussed in later sections.

It is important to emphasize that the flash tank is not an integral part of the condenser itself, even though it is often directly connected to it. The connection between these components is realized purely through piping, and their operation remains functionally independent. Efficient separation within the flash tank is nevertheless essential, as insufficient separation can lead to liquid carryover into the steam outlet or vapor presence in the condensate stream, both of which negatively affect the performance and reliability of downstream equipment such as condensers and pumps.

From a system perspective, flash tanks contribute to the overall efficiency of steam systems by enabling recovery and controlled routing of flash steam, as well as stabilization of hydraulic conditions in the condensate network \cite{salimi_technical-economic_2023}. In industrial practice, they are commonly employed wherever high-pressure condensate is expanded to lower pressure levels, making them a standard component of steam cycle auxiliary systems \cite{goodarzvand-chegini_energy_2023}.

The internal flow within a flash tank is characterized by complex multiphase behavior, including high-velocity inlet jets, strong turbulence, recirculation zones, and gravity-driven phase separation. These phenomena are strongly influenced by the geometry of the vessel and the configuration of inlet and outlet nozzles. As a result, simplified design assumptions may not fully capture the internal flow structure, which motivates the use of computational fluid dynamics (CFD) methods for detailed analysis of phase distribution, velocity fields, and separation efficiency.


\section{Vlastnosti páry a kondenzátu při náhlé změně tlaku}
This section summarizes the thermophysical phenomena associated with a sudden pressure reduction of condensate in systems connected to a condenser. The aim is to define the fundamental mechanisms leading to flashing, formation of a two-phase mixture, and the resulting flow conditions relevant for phase separation. Only those aspects that directly influence the internal flow behavior of the flash tank and its subsequent numerical modeling are considered.

\subsection{Termofyzikální vlastnosti páry a kondenzátu při náhlé změně tlaku}

Because the flash tank is connected to the condenser, it operates at one of the lowest pressure levels in the system, typically in the near-vacuum range. Under these conditions, incoming condensate undergoes a rapid pressure reduction, which strongly modifies its thermophysical state and promotes the formation of a vapor--liquid mixture. The subsequent behavior of this mixture is governed primarily by the large differences in density, specific volume, and compressibility between the liquid and vapor phases \cite{IAPWS_IF97, Ishii_Hibiki}.

A key parameter is the pronounced density difference between condensate and steam. Under near-vacuum conditions, the density ratio between the liquid and vapor phases typically reaches two to three orders of magnitude, while the specific volume of the vapor phase increases significantly \cite{IAPWS_IF97}. This disparity gives rise to strong buoyancy effects and creates favorable conditions for gravity-driven phase separation. At the same time, the rapid expansion of the vapor phase contributes to local flow acceleration and influences the velocity field established inside the vessel.

An important condition related to separation performance is the limitation of vapor-phase velocity with respect to droplet entrainment. In order to prevent liquid droplets from being carried into the vapor outlet, the upward steam velocity must remain below a critical value at which gravitational forces acting on the droplets are balanced by aerodynamic drag. In simplified form, this critical velocity may be expressed as \cite{Ishii_Hibiki}:

\begin{equation}
u_{\text{crit}} = \sqrt{\frac{4 g d (\rho_l - \rho_v)}{3 C_D \rho_v}}
\end{equation}

where $d$ is the droplet diameter, $\rho_l$ and $\rho_v$ are the densities of the liquid and vapor phases, $g$ is the gravitational acceleration, and $C_D$ is the drag coefficient. If this limit is exceeded, liquid droplets may be entrained into the vapor stream and transported toward the steam outlet, thereby reducing separation efficiency and potentially affecting downstream equipment.

The thermophysical properties described above define the basic conditions under which flash evaporation, phase redistribution, and subsequent internal flow behavior in the EPK take place. In particular, the large density contrast and the strong increase in vapor specific volume provide the physical basis for both phase separation and the development of characteristic internal flow structures. The mechanism of flash formation itself is discussed in the following subsection.

\subsection{Mechanismus náhlého vypaření (Flash) při poklesu tlaku}

Flashing is a thermodynamic process in which a liquid undergoes partial phase change due to a sudden pressure reduction below the saturation pressure corresponding to its initial state. In the case of a flash tank, this process is initiated when condensate is discharged from a higher-pressure system into a vessel operating at significantly lower pressure, typically near vacuum conditions. However, the exact location at which flashing occurs is not strictly confined to the vessel itself, as the pressure drop may already develop within the upstream piping or at the inlet restriction. Consequently, phase change can begin prior to entering the tank, and the inlet flow may already consist of a two-phase mixture. Under such conditions, the liquid contains excess specific enthalpy relative to the saturation state at the reduced pressure, and a portion of the liquid evaporates, forming flash steam \cite{Moran_Shapiro}. Assuming thermodynamic equilibrium, the resulting vapor fraction can be estimated from an energy balance as:

\begin{equation}
x = \frac{h_{in} - h_f}{h_{fg}}
\end{equation}

where $x$ is the vapor quality, $h_{in}$ is the specific enthalpy of the incoming fluid, $h_f$ is the saturated liquid enthalpy at the reduced pressure, and $h_{fg}$ is the latent heat of vaporization.

The extent of flashing is primarily governed by the magnitude of the pressure drop and the thermodynamic state of the incoming condensate. For saturated or near-saturated liquid, even a relatively small pressure reduction can result in significant vapor generation, whereas subcooled condensate requires additional energy to reach saturation before phase change occurs. Under near-vacuum conditions typical of condenser-connected flash tanks, the large difference in saturation properties leads to a substantial increase in specific volume of the vapor phase, which strongly influences the resulting flow field. In real systems, flashing may proceed under non-equilibrium conditions, where metastable liquid persists temporarily and vapor formation is affected by nucleation processes and local turbulence. However, for engineering analysis, the assumption of thermodynamic equilibrium is commonly adopted, as it provides a reasonable estimate of the vapor fraction and simplifies further modeling \cite{Ishii_Hibiki}.

The flashing process results in the formation of a vapor--liquid mixture characterized by a high volumetric fraction of vapor and dispersed liquid structures, such as droplets or continuous liquid regions. Due to the rapid expansion of the vapor phase, the flow is typically associated with high velocities, strong turbulence, and pronounced phase interaction. These effects govern the subsequent behavior within the flash tank, including phase distribution, droplet entrainment, and interaction with the vessel walls and free surface. From a modeling perspective, resolving the detailed phase change process is not essential for capturing these phenomena. Instead, the resulting two-phase mixture can be represented at the inlet by a prescribed vapor fraction, allowing the analysis to focus on the internal flow structure, surface behavior, and separation mechanisms under given operating conditions.

\subsection{Škrcení vody před vstupem do EPK}

The pressure reduction leading to flashing can be realized through a throttling process, in which the fluid passes through a localized flow restriction. In practical applications, a commonly used solution is a so-called dump tube, consisting of a perforated pipe with multiple small openings that distribute the flow and induce a pressure drop. The use of such a device allows the incoming high-momentum stream to be divided into multiple smaller jets, thereby reducing local velocities and limiting undesirable effects such as excessive turbulence, liquid entrainment, and direct impingement on the vessel walls. The throttling process is characterized by a rapid decrease in pressure with negligible heat exchange with the surroundings, and can therefore be approximated as isenthalpic, i.e., $h_{in} \approx h_{out}$ \cite{Moran_Shapiro}. As a consequence, the fluid retains its initial specific enthalpy while the saturation conditions shift to a lower pressure level, creating a thermodynamic imbalance that initiates flash evaporation.


\section{Vícefázové proudění uvnitř nádoby}

\subsection{Charakter dvoufázového proudění uvnitř expandéru}

After the pressure reduction and possible flash evaporation described in the previous section, the medium entering the EPK can no longer be treated as a single-phase fluid, but rather as a vapor–liquid mixture whose behavior is governed not only by thermodynamic state, but also by the internal hydrodynamics of the vessel. While Section 1.2 focused on the origin of this two-phase mixture and the mechanisms leading to its formation, the present subsection addresses its subsequent flow behavior inside the flash tank itself. In this context, the EPK must be understood not merely as a volume for pressure relief, but as a vessel in which the interaction of vapor, liquid, and the free surface determines the resulting flow structure, stability of the liquid region, and conditions for effective phase separation.

Uvnitř EPK dochází k charakteristickému prostorovému rozdělení obou fází, které je dáno především rozdílem jejich hustot, působením gravitace a uspořádáním samotné nádoby \cite{IshiiHibiki2011,HarunEtAl2019}. Parní fáze zaujímá převážně horní část expandéru a je odváděna parním výstupem, zatímco kapalná fáze se akumuluje ve spodní části nádoby a vytváří souvislý objem kondenzátu \cite{HarunEtAl2019}. Mezi oběma oblastmi se vytváří parokapalinové rozhraní, jehož tvar však za reálných provozních podmínek nemusí být rovinný ani stabilní \cite{IshiiHibiki2011,CarusoEtAl2016}. V blízkosti vstupní oblasti se navíc může vyskytovat lokálně neoddělená směs, zahrnující kapalné kapky unášené parní fází, pronikající kapalné proudy nebo oblasti s dosud neúplnou separací \cite{HarunEtAl2019}. Proudění uvnitř EPK proto nelze chápat jako dokonale ustálené a jednoznačně rozvrstvené, ale jako proměnný vícefázový systém, v němž se rozložení obou fází mění v prostoru i čase \cite{IshiiHibiki2011,CarusoEtAl2016}.

The internal flow field inside the EPK is governed primarily by the interaction of inlet momentum, gravity, the large density difference between the vapor and liquid phases, turbulence, and interfacial momentum exchange \cite{IshiiHibiki2011}. After entering the vessel, the fluid forms a jet or a system of jets whose penetration depth, direction, and dissipation depend on the inlet arrangement and local momentum flux; this inlet-driven motion redistributes both phases and may generate recirculation zones before gravity-controlled separation becomes dominant \cite{LeeEtAl2019}. At the same time, the strong density contrast between steam and condensate promotes stratification and buoyancy-induced rearrangement of the phases, whereas turbulent mixing acts in the opposite direction by increasing local interface deformation and delaying calm separation \cite{IshiiHibiki2011}. If the inlet configuration introduces a tangential velocity component, part of the inlet momentum may be transformed into rotational motion of the liquid, which can lead to swirl, vortex formation, and deformation of the free surface \cite{CarusoEtAl2016}. The resulting flow structure in the EPK is therefore determined by the combined action of inertial, gravitational, and turbulent effects rather than by thermodynamic considerations alone \cite{IshiiHibiki2011}.

Under these conditions, several characteristic flow phenomena may develop inside the EPK. In the inlet region, the incoming jet can penetrate into the vessel volume and generate local recirculation zones, which disturb the otherwise gravity-driven tendency toward phase stratification \cite{LeeEtAl2019,IshiiHibiki2011}. If the inlet orientation introduces a tangential component of velocity, the liquid phase may acquire a pronounced rotational motion, leading to swirl and, under sufficiently strong conditions, to vortex formation within the vessel \cite{CarusoEtAl2016,CristofanoEtAl2014}. Such rotating flow may subsequently deform the vapor--liquid interface and, in extreme cases, produce a paraboloidal free-surface shape, which is particularly undesirable in vessels where stable liquid level behavior is required \cite{CarusoEtAl2016,CristofanoEtAl2014}. At the same time, elevated local velocities and turbulent fluctuations can promote droplet entrainment into the vapor region, thereby worsening the separation process and increasing the risk of liquid carryover toward the steam outlet \cite{IshiiHibiki2011,LeeEtAl2019}. The internal flow in the EPK is therefore not defined only by bulk phase separation, but also by the presence and mutual interaction of recirculation, rotation, interface deformation, and entrainment phenomena \cite{IshiiHibiki2011}.

The practical importance of the above phenomena lies in their direct influence on the operational stability and separation performance of the EPK. Strong recirculation, excessive interface deformation, or persistent rotational motion may disturb the liquid level, worsen phase separation, and increase the risk of liquid carryover into the steam outlet, which can negatively affect downstream equipment and the overall reliability of the system \cite{IshiiHibiki2011,LeeEtAl2019}. For the sake of clarity, the mechanisms described in the previous paragraphs were presented primarily for the case of a single dominant inlet; in real applications, however, the EPK often receives several inflows that may operate simultaneously, so that their mutual interaction further modifies the internal flow field \cite{CarusoEtAl2016,SafikhaniEtAl2018}. Depending on the number, position, and orientation of the active inlets, the resulting flow may become more asymmetric, more rotational, or more unstable than in the simplified single-inlet case, which makes the prediction of the internal hydrodynamics considerably more difficult \cite{CarusoEtAl2016,SafikhaniEtAl2018}. For this reason, the character of parokapalinové proudění inside the EPK cannot be assessed reliably only from simplified assumptions, and a more detailed analysis of the effect of inlet arrangement is required in the subsequent parts of this thesis \cite{IshiiHibiki2011,LeeEtAl2019}.

\subsection{Typické režimy proudění a chování rozhraní uvnitř expandéru}

In the case of the EPK, the term \emph{flow regime} should not be understood primarily in the classical sense used for two-phase flow in pipes, where the emphasis is placed on one-dimensional phase arrangement along a confined channel. Instead, within a separation vessel such as the EPK, the relevant regime is defined by the overall spatial distribution of the vapor and liquid phases, the shape and stability of the vapor--liquid interface, and the character of the internal circulation induced by the inlet configuration and operating conditions \cite{IshiiHibiki2011,LeeEtAl2019}. The internal flow regime is therefore determined not only by the amount of each phase present, but also by the way in which momentum is introduced into the vessel and subsequently redistributed by gravity, turbulence, and interfacial interaction \cite{IshiiHibiki2011}. For this reason, the description of flow regimes in the EPK must focus on characteristic internal flow states of the vessel rather than on conventional pipe-flow classifications \cite{IshiiHibiki2011}.

The simplest and, from the point of view of phase separation, also the most favorable internal regime is a predominantly stratified configuration. In this state, the liquid phase forms a continuous hold-up in the lower part of the vessel, while the vapor phase occupies the upper region and is discharged through the steam outlet; the interface between the two phases remains relatively smooth and only weakly disturbed \cite{IshiiHibiki2011,HarunEtAl2019}. Under such conditions, gravity-driven separation dominates over inertial mixing, and the internal circulation is limited to weak local motion that does not substantially disrupt the phase boundary \cite{IshiiHibiki2011}. Although this regime represents an idealized reference rather than a fully permanent condition in real operation, it provides a useful baseline for evaluating more disturbed flow states that arise when inlet momentum, turbulence, or rotational motion become significant \cite{HarunEtAl2019,LeeEtAl2019}.

When the momentum of the incoming stream becomes sufficiently high, the internal flow departs from the nearly stratified regime and enters an inlet-disturbed state. In this regime, the entering fluid forms a penetrating jet or a locally mixed inlet zone, in which the two phases are not immediately separated and the interface may become strongly deformed by direct momentum transfer \cite{LeeEtAl2019,HarunEtAl2019}. Depending on the position and orientation of the inlet, this disturbance may generate local waves, transient submergence of the free surface, or recirculation cells that transport liquid away from the lower hold-up region and delay the re-establishment of calm stratification \cite{LeeEtAl2019,IshiiHibiki2011}. The resulting regime is therefore controlled not only by the magnitude of the inlet flow rate, but also by the manner in which the momentum is introduced into the vessel volume, which makes the internal hydrodynamic response highly sensitive to inlet arrangement \cite{LeeEtAl2019,IshiiHibiki2011}.

A further characteristic regime may arise when the inlet configuration introduces a significant tangential component of momentum into the liquid phase. In such a case, the internal flow can become swirl-dominated, with the liquid rotating around the vessel axis and forming a circumferential velocity field that differs fundamentally from the nearly stratified or only locally disturbed states described above \cite{CarusoEtAl2016,CristofanoEtAl2014}. As the rotational intensity increases, the free surface may no longer remain approximately horizontal, but instead becomes progressively deformed into a curved shape, which can evolve toward a paraboloidal interface in the mean sense \cite{CarusoEtAl2016,CristofanoEtAl2014}. At the same time, the swirling motion may generate vortex structures and secondary recirculation patterns that redistribute liquid within the vessel and may locally amplify interface instability \cite{IshiiHibiki2011,CarusoEtAl2016}. This regime is particularly important in the case of the EPK, because even if it may support partial phase segregation in some regions, it can simultaneously create highly unfavorable conditions for stable level behavior and predictable internal flow organization \cite{CristofanoEtAl2014,IshiiHibiki2011}.

In real operating conditions, the internal flow inside the EPK does not remain permanently in one of the idealized regimes described above, but may transition between them depending on the active combination of inlets, the instantaneous flow rates, and the liquid level inside the vessel \cite{IshiiHibiki2011,LeeEtAl2019}. A regime that is predominantly stratified under one load case may become locally disturbed or even swirl-dominated when the inlet momentum increases or when several inlets interact simultaneously \cite{SafikhaniEtAl2018,LeeEtAl2019}. These transitions are important because they directly affect the stability of the vapor--liquid interface, the redistribution of liquid within the vessel, and the conditions under which droplets or larger liquid structures may be transported into the vapor region \cite{IshiiHibiki2011,HarunEtAl2019}. For this reason, the identification of internal flow regimes in the EPK cannot be regarded as an end in itself, but rather as a basis for understanding the associated phenomena of entrainment, recirculation, and surface instability, which are discussed in the following subsection \cite{IshiiHibiki2011}.

\subsection{Unášení kapaliny a nestabilita volné hladiny}

The flow regimes described in the previous subsection are not important only as different hydrodynamic states of the EPK, but primarily because they directly influence its separation performance and operational stability. Even when the general distribution of vapor and liquid inside the vessel is known, disturbed flow may still lead to undesirable consequences such as liquid entrainment into the vapor region or instability of the free surface, both of which can degrade the intended function of the expandor \cite{IshiiHibiki2011,LeeEtAl2019}. From an engineering point of view, these effects are particularly important because they determine whether the EPK behaves as a stable separation vessel or as a source of fluctuating internal flow conditions that may negatively affect downstream operation \cite{IshiiHibiki2011}.

One of the most important adverse consequences of disturbed internal flow is liquid entrainment into the vapor region. In the context of the EPK, entrainment refers to the transport of liquid from the main condensate hold-up or from the deformed free surface into the steam-dominated part of the vessel, where it may subsequently be carried toward the vapor outlet \cite{IshiiHibiki2011}. This process can be promoted by elevated local vapor velocities, strong interface deformation, wave breaking, or by the interaction of rotational flow with the free surface, all of which increase the likelihood that droplets or larger liquid structures will be detached from the liquid region and transported upward \cite{LeeEtAl2019,CarusoEtAl2016}. If such liquid leaves the vessel together with the separated steam, the result is carryover, which represents a direct reduction in separation quality and may negatively affect the operation of downstream equipment \cite{IshiiHibiki2011,HarunEtAl2019}.

In addition to entrainment, disturbed internal flow may also lead to instability of the free surface, which is of particular importance in the case of the EPK. Instead of remaining approximately steady, the vapor--liquid interface may exhibit waves, oscillatory motion, local depressions, or a sustained curved shape caused by rotational flow, and under sufficiently strong swirl the mean surface may approach a paraboloidal form \cite{CarusoEtAl2016,CristofanoEtAl2014}. Such behavior is undesirable because it reduces the stability of the liquid region, complicates the interpretation of level-related behavior, and indicates that the internal flow is no longer governed primarily by calm gravity-driven separation \cite{IshiiHibiki2011}. For this reason, the assessment of free-surface stability represents an essential part of evaluating the hydrodynamic performance of the EPK and forms an important basis for the later CFD analysis presented in this thesis \cite{IshiiHibiki2011,LeeEtAl2019}.